\documentclass[11pt]{article}
\usepackage[margin=1in]{geometry}
\usepackage{xcolor}
\usepackage{listings}
\usepackage{hyperref}
\usepackage{amsmath}
\usepackage{enumitem}
\hypersetup{colorlinks=true, linkcolor=blue!50!black, urlcolor=blue!50!black}

\definecolor{codebg}{gray}{0.96}
\lstset{
  basicstyle=\ttfamily\small, backgroundcolor=\color{codebg},
  breaklines=true, columns=fullflexible, frame=single, framesep=4pt,
  xleftmargin=4pt, xrightmargin=4pt, aboveskip=6pt, belowskip=6pt,
  keepspaces=true, showstringspaces=false
}
\newcommand{\code}[1]{\texttt{#1}}

\title{\vspace{-1cm}A Student's Guide to the $^{16}$C$(p,d)^{15}$C Analysis\\[2pt]
       \large (a1975, AT-TPC active target, genfit reconstruction)}
\author{}
\date{}

\begin{document}
\maketitle
\vspace{-1.2cm}

\section{What you are measuring}
The a1975 experiment is a $^{16}$C beam ($\sim$192~MeV, $\sim$12~MeV/u) stopped in an
\textbf{H$_2$ gas active target} at 1~bar. The gas is \emph{both} the proton target and the
tracking medium. You will analyze the \textbf{neutron pickup} reaction
\[
  ^{16}\mathrm{C} + p \;\rightarrow\; ^{15}\mathrm{C}^{*} + d,
\]
i.e.\ $^{16}$C$(p,d)^{15}$C, by detecting the outgoing \textbf{deuteron}. From the deuteron's
kinetic energy $T_d$ and lab angle $\theta_{\mathrm{lab}}$ you reconstruct the $^{15}$C
excitation energy $E_x$ via two-body relativistic kinematics, and you study which $^{15}$C
states are populated and how strongly (the angular distribution).

\medskip
\noindent\textbf{Goal of this guide:} starting from the reconstructed data, produce a clean
$^{15}$C excitation-energy spectrum and its angular dependence, and understand the two
systematics that matter: \emph{triton contamination} and the \emph{$E_x$-vs-$\theta_{cm}$ drift}.

\section{Setup}
All commands below are run from the (p,d) macro directory after sourcing the environment:
\begin{lstlisting}[language=bash]
cd ~/fair_install/ATTPCROOTv2-OpenKF/macro/Unpack_HDF5/a1975/UKF
source ~/fair_install/ATTPCROOTv2-OpenKF/build/config.sh
\end{lstlisting}
The raw data and reconstructed files live on the \code{F:} drive (mount it in WSL with
\code{sudo mount -t drvfs F: /mnt/f} if needed):
\begin{itemize}[nosep]
  \item \code{/mnt/f/a1975/h5/}\hfill raw HDF5 (one file per run, \code{run\_0106}--\code{run\_0189})
  \item \code{/mnt/f/a1975/reco/}\hfill reconstructed \code{*\_reco.root} (hits + pattern/tracks) + \code{*\_FRIB.root} (ion chamber)
  \item \code{/mnt/f/a1975/reco\_pd/}\hfill your (p,d) fit outputs go here
\end{itemize}
The analysis macros are in \code{pp/}; the interactive ones need an X display
(they work over WSLg; the static ones write a PNG you open with \code{explorer.exe}).

\section{The pipeline at a glance}
\begin{center}
\code{raw HDF5} $\to$ \fbox{unpack + PSA + space-charge + PRA} $\to$ \code{*\_reco.root}
$\to$ \fbox{deuteron genfit fit} $\to$ \code{*\_genfitter\_pd.root}
$\to$ \fbox{cache} $\to$ \code{pd\_kin.root} $\to$ \fbox{plots}
\end{center}
\textbf{Key idea:} pattern recognition (PRA, finding the tracks) does not depend on what
particle you assume, so the \code{*\_reco.root} files are shared by all channels. For (p,d) you
only need to \emph{fit} the existing tracks with the \textbf{deuteron} hypothesis. You do
\emph{not} re-run the reconstruction.

\section{Step 1 --- Fit the deuterons}
\code{fitGenfitter\_a1975.C} reads a \code{*\_reco.root}, keeps only tracks inside the deuteron
PID band (\code{pid/deuteron\_band.json}), fits each as a deuteron with genfit, and writes the
fitted kinematics. Its defaults are already the deuteron hypothesis, so for one run:
\begin{lstlisting}[language=bash]
root -l -b -q 'pipeline/fitGenfitter_a1975.C("run_0106", -1, "/mnt/f/a1975/reco/", "_pd", "/mnt/f/a1975/reco_pd/")'
\end{lstlisting}
The arguments are \code{(run, nEvents=-1 for all, inputDir, outputSuffix, outputDir)}. The output
is \code{/mnt/f/a1975/reco\_pd/run\_0106\_genfitter\_pd.root}. Run the whole experiment on 4 cores:
\begin{lstlisting}[language=bash]
ls /mnt/f/a1975/reco/run_*_reco.root | grep -oE 'run_[0-9]+' | sort -u > /tmp/pd_runs.txt
cat /tmp/pd_runs.txt | xargs -P4 -I{} \
  root -l -b -q "pipeline/fitGenfitter_a1975.C(\"{}\", -1, \"/mnt/f/a1975/reco/\", \"_pd\", \"/mnt/f/a1975/reco_pd/\")"
\end{lstlisting}
A deuteron is rare ($\sim$1.5\% of tracks), so most events show \code{0/1 tracks fitted} --- that
is the PID gate doing its job, not an error. Expect $\sim$30{,}000 deuterons over the 84 runs.

\section{Step 2 --- Cache the kinematics}
The fit files are large; cache just the per-deuteron numbers you need into a flat ntuple. The
cache also pulls the \textbf{ion-chamber (IC)} value from the \code{*\_FRIB.root} files
(the IC tags good $^{16}$C beam particles).
\begin{lstlisting}[language=bash]
cat /tmp/pd_runs.txt | xargs -P4 -I{} root -b -q 'pp/cache_pd_run.C("{}","/tmp/pdkin_{}.root")'
hadd -f /tmp/pd_kin.root /tmp/pdkin_run_01*.root
\end{lstlisting}
This makes \code{/tmp/pd\_kin.root}, an ntuple \code{pk} with branches
\code{ke, theta, vz, chi2ndf, ic, run}. All later plots read this single file (fast).

\section{Step 3 --- The excitation-energy spectrum}
From $T_d$ and $\theta_{\mathrm{lab}}$ the macros compute $E_x(^{15}\mathrm{C})$ with relativistic
two-body kinematics. Make the 4-panel overview:
\begin{lstlisting}[language=bash]
root -l -b -q 'pp/plot_pd.C'                          # -> /tmp/pd_spectrum.png
explorer.exe "$(wslpath -w /tmp/pd_spectrum.png)"
\end{lstlisting}
You should see (i) the $E_x$ spectrum with a sharp $^{15}$C ground-state peak near 0 and excited
states ($\sim$3.10, 4.66~MeV); (ii) the deuteron $T_d$-vs-$\theta_{\mathrm{lab}}$ band lying on the
analytic kinematic line; (iii) $E_x$ vs $\theta_{cm}$; (iv) vertex-$z$ vs $E_x$.

\paragraph{Cuts.} \code{plot\_pd.C} applies a $\chi^2/\mathrm{ndf}$ cut and the IC beam gate
($950$--$1350$) by default. These remove badly-fit and non-beam events.

\section{Step 4 --- Angular dependence}
A discrete state should appear at the same $E_x$ at every angle, and its population (cross
section) changes with angle. Plot $E_x$ in $\theta_{cm}$ slices:
\begin{lstlisting}[language=bash]
root -l -b -q 'pp/plot_pd_angbins.C(10,5,12)'         # start 10 deg, 5 deg wide, 12 bins
explorer.exe "$(wslpath -w /tmp/pd_angbins.png)"
\end{lstlisting}
Arguments are \code{(thcmStart, thcmWidth, nbins, \dots)}. The ground-state resolution is best at
forward angles ($\sigma\approx0.15$~MeV at $20$--$30^\circ$) and degrades toward back angles.

\section{Step 5 --- Correct the $E_x$-vs-$\theta_{cm}$ drift (important!)}
The reconstructed ground state is \emph{not} flat in $\theta_{cm}$: its centroid drifts by
$\sim$0.9~MeV across the angular range. This is a kinematic/reconstruction systematic, and it
\textbf{smears the integrated spectrum}. The fix measures the g.s.\ centroid versus $\theta_{cm}$
and subtracts it event-by-event:
\begin{lstlisting}[language=bash]
root -l -b -q 'pp/excorr_pd.C'                         # -> /tmp/pd_excorr.png
\end{lstlisting}
The left panel shows the measured drift (and the curve subtracted); the right panel shows the
integrated $^{15}$C spectrum before and after. The ground-state FWHM improves from
$\sim$1.7~MeV to $\sim$0.43~MeV. \emph{This is the single most important correction for the
integrated spectrum.} The same correction is built into \code{plot\_pd\_angbins.C}
(argument \code{driftCorr}, on by default) and into the interactive explorer (a checkbox).

\section{Step 6 --- Beware triton contamination}
$^{16}$C$(p,t)^{14}$C produces \textbf{tritons}. The deuteron PID gate is loose enough to admit
some of them, and genfit fits a triton \emph{as a deuteron} with a perfectly good $\chi^2$
(so you cannot remove tritons with a $\chi^2$ or IC cut --- we checked). They show up as a faint
second band above the deuteron locus in $T_d$ vs $\theta_{\mathrm{lab}}$. Look at the PID plane:
\begin{lstlisting}[language=bash]
root -l -b -q 'pp/pid_pd.C(106,113)'                   # sqrt(dE/dx) vs Brho + the gate
explorer.exe "$(wslpath -w /tmp/pid_pd.png)"
\end{lstlisting}
The proton, deuteron and triton bands are separated in the $\sqrt{dE/dx}$--$B\rho$ plane; the red
polygon is the current (loose) deuteron gate. \textbf{The fix is to tighten the gate} so the
triton band is excluded --- because the species separation must come from the gate, not from the
fit. Use the interactive editor (draw a polygon, save it to JSON, then re-run Step~1 with the new
gate):
\begin{lstlisting}[language=bash]
root -l 'pp/pid_gate_edit.C'      # Draw new gate -> click vertices, double-click to close -> Save JSON
\end{lstlisting}

\section{Step 7 --- Interactive exploration}
For tuning cuts and binning live (needs a display):
\begin{lstlisting}[language=bash]
root -l 'pp/explore_pd.C'
\end{lstlisting}
Controls: beam energy, $\chi^2/\mathrm{ndf}$, $\theta$ window, IC window, \textbf{vertex-$z$
bounds}, $E_x$ binning, plus check-boxes for the \textbf{drift correction} and the $^{15}$C
reference lines. Adjust and press \code{Redraw}; \code{Save PNG} to export.

\section{Physics checklist}
\begin{itemize}[nosep]
  \item The $^{15}$C ground state ($1/2^+$) should sit at $E_x=0$ after calibration; the first
        excited state is at 0.740~MeV ($5/2^+$, on the g.s.\ shoulder).
  \item Prominent strength near 3.10 and 4.66~MeV; broad structure at 6--8~MeV.
  \item The deuteron $T_d(\theta)$ band must follow the analytic ground-state kinematic line --- if
        it does not, suspect a wrong beam energy, $B$-field sign, or frame convention.
\end{itemize}

\section{Troubleshooting}
\begin{itemize}[nosep]
  \item \emph{``Macro is slow / hangs.''} The PID macros only need the \code{AtPIDEvent} branch;
        they already disable the big track branch. The IC gate reads the large FRIB files and is
        therefore slower --- it is opt-in in \code{pid\_pd.C}.
  \item \emph{``\code{/tmp/pd\_kin.root} missing.''} \code{/tmp} clears on reboot --- rebuild the
        cache (Step~2).
  \item \emph{``Gate file not found.''} Run from the \code{UKF} directory so \code{pid/...json}
        resolves.
  \item \emph{``Interactive window does not open.''} The explorer/editor need X (WSLg); the
        \code{-b} flag means batch --- drop it for the GUIs.
\end{itemize}

\end{document}
